\documentclass[12pt,a4paper]{article}

\usepackage[spanish,es-nodecimaldot]{babel}
\usepackage[utf8]{inputenc}
\usepackage[T1]{fontenc}
\usepackage{geometry}
\usepackage{graphicx}
\usepackage{float}
\usepackage{caption}
\usepackage{hyperref}
\usepackage{enumitem}
\geometry{margin=2.5cm}
\graphicspath{{../../evidencias/}}
\hypersetup{colorlinks=true,linkcolor=blue,urlcolor=blue}

\title{Informe técnico\\Taller Integrado de Servicios de Red}
\author{Luis Jaramillo \\ Giovanny Toledo \\ Maximiliano Rivas\\Grupo 3}
\date{\today}

\begin{document}
\maketitle
\tableofcontents
\clearpage

\section{Introducción}
% Reemplazar esta nota por el alcance y contexto confirmados en la pauta vigente.
Pendiente: describir el objetivo del taller y el alcance efectivamente trabajado.

\section{Objetivos}
\subsection{Objetivo general}
Pendiente: transcribir o resumir el objetivo vigente de la actividad.
\subsection{Objetivos específicos}
Pendiente: enumerar los objetivos confirmados en la pauta.

\section{Arquitectura y decisiones}
Pendiente: incluir un diagrama y explicar los componentes, sus relaciones y las decisiones adoptadas. Distinguir servicios planificados de servicios verificados.

\section{Entorno y línea base}
Pendiente: documentar sistema operativo, recursos y estado inicial medidos en el VPS. No incluir contraseñas ni llaves.

\section{Implementación y pruebas}
En cada subsección explicar qué se hizo, cómo se configuró, por qué se tomaron las decisiones relevantes, qué prueba se ejecutó y qué resultado obtuvo. Las figuras deben llevar una explicación y acreditar un punto concreto.

\subsection{Preparación del VPS}
Pendiente: configuración y pruebas propias del Grupo 3.

\subsection{Usuarios y directorios}
Pendiente: configuración, permisos y comprobaciones.

\subsection{DNS y base de datos}
Pendiente: configuración, decisiones y consultas de validación.

\subsection{Administración DNS}
Pendiente: acceso, cambios y prueba observable.

\subsection{Servidor web y CMS}
Pendiente: documentar cada sitio y CMS instalado, su versión y las pruebas correspondientes.

\subsection{FTP}
Pendiente: configuración, aislamiento y pruebas de transferencia.

\subsection{Correo}
Pendiente: configuración y pruebas de envío y recepción requeridas.

\subsection{Firewall y SELinux}
Pendiente: reglas, estado, ajustes y pruebas de conectividad.

\subsection{Aplicación de sockets}
Pendiente: describir el programa del equipo, protocolo, puerto y prueba desde otro equipo.

\subsection{Verificación integrada}
Pendiente: documentar pruebas de integración y sus resultados.

\section{Problemas y soluciones}
Pendiente: registrar solo incidentes que ocurrieron, su causa comprobada, la solución aplicada y su validación.

\section{Conclusiones}
Pendiente: resumir el estado real de los objetivos y las tareas que queden parciales o pendientes.

\end{document}
